\documentclass[11pt,a4paper]{article}

\usepackage[T1]{fontenc}
\usepackage[utf8]{inputenc}
\usepackage{amsmath,amssymb}
\usepackage{booktabs}
\usepackage{listings}
\usepackage{url}
\usepackage[margin=1in]{geometry}
\usepackage[hidelinks]{hyperref}

\lstset{basicstyle=\ttfamily\small,columns=fullflexible,breaklines=true}

\newcommand{\code}[1]{\texttt{#1}}
\newcommand{\us}{\,$\mu$s}

% Paragraphs dense with \code names would otherwise overrun the margin.
\setlength{\emergencystretch}{3em}

\title{GMP on the CPU or CGBN on the GPU?\\[4pt]
  \large Measurements from gmp-ecpp}

\author{Stefan Teleman \code{<stefan.teleman@gmail.com>}\\[4pt]Claude Opus 5.5 (Anthropic)}

\date{November 17, 2026}

\begin{document}
\maketitle

\begin{abstract}
This report asks whether the multiple-precision arithmetic that gmp-ecpp
does on the host with GMP should move to the GPU with CGBN, in its ECPP
prover (\code{libgmpecpp}, \code{gmpecpp-mpi}) and in its ECM factorizer
(\code{libgmpecmfft}, \code{gmpecmfft-mpi}). Three experiments answer it:
the latency and throughput of a single operation, where each program's
host time actually goes, and end-to-end batch proving over MPI with and
without the CGBN backend, across several layouts and seeds. Replacing GMP
calls one for one with CGBN is $4$ to $172\times$ slower; batched, CGBN is
$12$ to $124\times$ faster, and the work that batches has already moved.
End to end, the CGBN backend's benefit falls as the host gets threads,
from $1.83\times$ to about parity, while it consistently removes
pathological proofs. A default of Open MPI's process binding can fake a
$3.6\times$ advantage. Every number comes from a committed harness
(\code{make bench-cgbn-latency}, \code{bench-gmp-profile},
\code{bench-mpi-cgbn}).
\end{abstract}

\tableofcontents

\section{Summary}

\begin{itemize}
  \item \textbf{One-for-one replacement of GMP calls with CGBN is slower,
    by $4\times$ to $172\times$.} A single CGBN instance runs on 8--32
    threads of one SM. Every call pays a fixed round trip of about
    26\us: copies in, a launch, a synchronization, a copy out. GMP
    finishes a 512-bit modular multiplication in 0.15\us.
  \item \textbf{Batched, CGBN is $12\times$ to $124\times$ faster than
    GMP.} That takes thousands of independent operations of one width in
    a single launch.
  \item \textbf{The batchable host work has already moved; the rest
    cannot.}
    \begin{itemize}
      \item In the factorizer, host GMP is at most 0.06\% of host time,
        and above 16352 bits CGBN cannot hold the numbers at all.
      \item In the prover, the batchable part is the elliptic curve method
        (ECM), and \code{-C} already runs it on CGBN. What GMP still does
        is mostly affine point arithmetic (gcd-bound), in work that is
        small and sequential: the host's deliberate first attempt at each
        factorization, and the downrun's point checks.
        Experiment~1 (\S\ref{sec:exp1}) rules out one-call-at-a-time
        offload of either.
    \end{itemize}
  \item \textbf{End to end, the CGBN backend's benefit depends on the
    host's thread budget.}
    \begin{itemize}
      \item It is $1.83\times$ with one worker of 8 ECM threads,
        $1.37\times$ with 16, and $1.20\times$ with 32.
      \item With 4 or 8 workers sharing one GPU it is roughly even
        ($0.98\times$, $1.16\times$).
      \item In every configuration it cuts variance: one worker's
        18-proof batch ranged from 39 to 57\,s on the host and from 24
        to 28\,s with CGBN.
    \end{itemize}
  \item \textbf{Process binding can fake the result.} \code{mpirun -n 2}
    binds the worker to one core by default. All its ECM threads then
    share that core, the host path slows $4\times$, and a spurious
    $3.6\times$ ``GPU advantage'' appears (\S\ref{sec:binding}).
\end{itemize}

\section{System under test}

\begin{center}
\begin{tabular}{ll}
\toprule
CPU & Intel Core i9-13900K: 8 P-cores $\times$ 2 threads (CPUs 0--15) \\
    & + 16 E-cores (CPUs 16--31), up to 5.8\,GHz \\
GPU & NVIDIA RTX PRO 4500 Blackwell (sm\_120), 32\,GB, driver 595.71.05 \\
CUDA & 13.2; CGBN (NVlabs) \\
GMP & 6.3.0 (Fedora 41) \\
MPI & Open MPI 5.0.5, rebuilt CUDA-aware (irrelevant here: \\
    & all messages are small host messages) \\
OS & Fedora 41, Linux 6.17.10 \\
Compilers & GCC 14 (host); \code{nvcc -O3 -arch=sm\_120} \\
\bottomrule
\end{tabular}
\end{center}

\section{Background: what runs where}

\paragraph{CGBN} is NVlabs' fixed-width multiple-precision library. An
instance of \code{BITS} bits is spread over \code{TPI} threads (8, 16 or
32; cooperative groups, warp-synchronous), each holding $\code{BITS}/(32
\cdot \code{TPI})$ 32-bit limbs. Arithmetic is done cooperatively across
those threads; modular multiplication is Montgomery's. Throughput comes
only from running many independent instances per launch. In this project
the widest instance is 16384 bits, because at $\code{TPI} = 32$ the
kernels already use 236--248 of 255 registers.

\paragraph{The prover} (Atkin--Morain ECPP) is mostly host GMP: \code{mpz}
for the number theory, \code{mpf} for the class polynomials, polynomial
root finding mod $p$, and the elliptic curve method, which it uses to
factor candidate curve orders. Its host ECM (\code{LenstrasECM}) works in
affine Weierstrass coordinates, so every point addition needs a modular
inverse, that is, an extended gcd. With \code{-T N} it runs $N$ threads
(\code{LenstrasECM\_mt}):
\begin{itemize}
  \item each thread takes curves $\mathit{id}, \mathit{id}+N, \ldots$
    from $\max(N, 20)$ curves per bound, and runs its own sequence of
    bounds without a barrier;
  \item the first thread to find a factor sets an atomic flag, and the
    others stop.
\end{itemize}
With \code{-C}, the prover keeps the host attempt first and escalates to
the GPU only when the cofactor is at least 192 bits. The GPU then runs 4096
Montgomery-curve instances (no inversions) at once, stage one and stage
two, with $B_1$ clamped to $[10^3, 10^5]$. An earlier version escalated
every failed host attempt and made proofs $2\times$ slower: most calls are
meant to fail fast, because they are rejecting candidate curve orders.

\paragraph{The factorizer} runs ECM entirely on the GPU, on CGBN up to
16352 bits and on cuFFT arithmetic above. The host does curve setup, moves
a few values between engines, and takes one gcd per curve at the end.

\paragraph{\code{gmpecpp-mpi}} distributes whole numbers: rank 0 hands out
indices on request, and each worker calls \code{gmpecpp\_prove()}. Every
proof reseeds the prover's generator from the configured seed, so the
order in which ranks take numbers does not affect any proof; \code{-s}
sets that seed.

\section{Experiment 1: one operation, one instance, or 4096}
\label{sec:exp1}

\paragraph{Method} (\code{bench/cgbn-latency.cu}). For 512, 1024 and 2048
bits (CGBN \code{TPI} = 8, 16, 32): modular exponentiation $a^e \bmod n$
with a full-width exponent, the core of Rabin--Miller and of square roots
mod $p$, and modular multiplication $ab \bmod n$ (\code{cgbn\_mul\_wide}
+ \code{cgbn\_rem\_wide}).
\begin{itemize}
  \item \textbf{GMP:} one thread, \code{mpz\_powm}, and \code{mpz\_mul} +
    \code{mpz\_mod}.
  \item \textbf{CGBN, one:} a full round trip for one instance: three
    \code{cudaMemcpy} host to device, a \code{<<<1, TPI>>>} launch, and a
    \code{cudaMemcpy} device to host, which synchronizes. This is the cost
    of replacing one GMP call.
  \item \textbf{CGBN, batch:} 4096 instances in one launch (128 threads
    per block), plus the result copy, divided by 4096. Contexts use
    \code{cgbn\_no\_checks}.
  \item One result per width was compared with GMP; all matched.
\end{itemize}

\begin{table}[ht]
\centering
\begin{tabular}{lrrr}
\toprule
Operation & GMP & CGBN, one instance & CGBN, 4096 per launch \\
\midrule
$a^e \bmod n$, 512 bits  & 49\us   & 1439\us{} ($29\times$ slower) & 0.83\us/op ($59\times$ faster) \\
$a^e \bmod n$, 1024 bits & 321\us  & 4774\us{} ($15\times$ slower) & 3.14\us/op ($102\times$ faster) \\
$a^e \bmod n$, 2048 bits & 2157\us & 9210\us{} ($4.3\times$ slower) & 17.4\us/op ($124\times$ faster) \\
$ab \bmod n$, 512 bits   & 0.15\us & 26.2\us{} ($172\times$ slower) & 0.013\us/op ($12\times$ faster) \\
$ab \bmod n$, 1024 bits  & 0.43\us & 26.2\us{} ($61\times$ slower) & 0.019\us/op ($22\times$ faster) \\
$ab \bmod n$, 2048 bits  & 1.42\us & 26.3\us{} ($19\times$ slower) & 0.036\us/op ($40\times$ faster) \\
\bottomrule
\end{tabular}
\caption{One GMP operation against one CGBN instance and against 4096.}
\label{tab:lat}
\end{table}

\paragraph{Reading it.}
\begin{itemize}
  \item \textbf{The fixed cost.} The multiplication rows are flat at about
    26\us{} for one instance at every width. That is the round trip
    (pageable \code{cudaMemcpy}, launch, synchronization), not the
    arithmetic. Any design that hands single operations to the GPU pays it
    every time; GMP's 512-bit multiplication is 170 times shorter.
  \item \textbf{The parallelism.} A single exponentiation leaves the GPU
    almost idle: one instance is one warp, or a quarter of one, on one SM
    of many, executing a serial square-and-multiply chain. Against it is
    one host core at up to 5.8\,GHz running GMP's hand-tuned x86-64
    assembly. The gap narrows with width ($29\times$ at 512 bits,
    $4.3\times$ at 2048) because more threads per instance share the work.
    It does not close.
  \item \textbf{The batch.} With 4096 instances per launch the device is
    full and the per-operation cost collapses. The advantage grows with
    width for exponentiation ($59 \to 124\times$), where more arithmetic
    sits behind each byte copied.
  \item \textbf{Not measured:} launch-latency tuning (CUDA graphs, pinned
    memory, persistent kernels) would shave part of the 26\us. No such
    technique removes the gap between one warp and the whole GPU.
\end{itemize}
CGBN is worth using only where an algorithm exposes thousands of
independent same-width operations at once. ECM curves do, and so would,
for example, primality tests of many candidates at once. A sequential
chain does not: a downrun step, a gcd, or Cornacchia on one discriminant.

\section{Experiment 2: where the host time goes}
\label{sec:exp2}

\paragraph{Method} (\code{bench/gmp-profile.sh}). \code{perf record -e
cpu-clock -F 1999} on each run, then \code{perf report -{}-sort dso}; for
the prover also \code{-{}-sort dso,sym}, with \code{libgmp}'s symbols grouped
into gcd and division (\code{*gcd*}, \code{*hgcd*}, \code{div2},
\code{*divrem*}, \code{*div\_qr*}, \code{*bdiv*}, \code{*invert*}),
\code{mpf\_*}, and the rest. The \code{libgmp} share bounds what any GPU
offload of GMP could save. A thread blocked in a CUDA synchronization
spins inside \code{libcuda}, so that share is time spent waiting on the
GPU.

On this hybrid CPU the default \code{cycles} event is two PMU events,
\code{cpu\_core/cycles/} and \code{cpu\_atom/cycles/}. \code{perf report}
prints one table per event, and a thread that migrates between core types
is split across both. \code{cpu-clock} is a software timer: it counts
time, uniformly, in one table.

\begin{table}[ht]
\centering
\begin{tabular}{lrrrrr}
\toprule
Run & Wall & \code{libgmp} & gcd, division & \code{mpf\_*} & \code{libcuda} \\
\midrule
Prover, 140 digits, host only, 1 thread & 67.6\,s & 92.2\% & 76.0\% & 0.0\% & --- \\
Prover, same, \code{-C}                 & 2.9\,s  & 64.8\% & 52.0\% & 0.1\% & 23.6\% \\
\code{gmpecmfft}, 2000 bits (CGBN)      &         & 0.04\% &        &       & 90\% \\
\code{gmpecmfft}, 20000 bits (cuFFT)    &         & 0.06\% &        &       & 89\% \\
\bottomrule
\end{tabular}
\caption{Share of host time by library (\code{perf}, \code{cpu-clock} at
  1999\,Hz). The prover runs seed 1; \code{gmpecmfft} runs 2952 and 64
  curves, $B_1 = 10^4$, $B_2 = 10^6$.}
\label{tab:prof}
\end{table}

\paragraph{Reading it.}
\begin{itemize}
  \item \textbf{The factorizer} is essentially pure GPU. Its host GMP is at
    most 0.06\% of host time, and the host thread spends about 90\% waiting
    on the device (most of the rest is \code{[vdso]}, the polling loop's
    clock reads). No GMP replacement can show up here; above 16352 bits
    there is nothing CGBN could do anyway.
  \item \textbf{The prover, host only,} is dominated by gcd and division:
    \code{\_\_gmpn\_hgcd2} about 39\% and \code{div2} about 22\% of
    samples. The prover's point arithmetic (\code{multiply}) is affine, with
    an inversion at every point addition, and it serves the host ECM and
    the downrun's point checks alike. The ECM share is precisely what
    \code{-C} replaces: CGBN does it with Montgomery's x-only curves, which
    need no inversion, across 4096 curves at once. Hence $67.6 \to 2.9$\,s,
    $23\times$, on this proof. A frame-pointer call graph cannot divide
    \code{multiply}'s time between its two callers: GMP's assembly keeps no
    frame pointer, so unwinding stops early.
  \item \textbf{The prover, \code{-C}:} 65\% of the time remains in GMP,
    still mostly gcd and division (52\%); \code{mpf} is 0.1\% at the
    default precision, and the \code{libcuda} 24\% is the host waiting on
    the GPU's ECM. The affine arithmetic that remains belongs to small,
    sequential work: the host's first attempt at every factorization, kept
    on the host because most attempts are meant to fail quickly
    (escalating all of them made proofs $2\times$ slower), and the
    downrun's point checks, a short chain per step. Neither offers
    thousands of independent same-width operations. Batching the host
    attempts into a small GPU ECM is conceivable but untested, and it would
    have to beat a fixed per-launch cost that most of these attempts,
    failing in milliseconds, cannot amortize.
\end{itemize}

\section{Experiment 3: batch proofs over MPI}
\label{sec:exp3}

\paragraph{Method} (\code{bench/mpi-cgbn.sh}, \code{make bench-mpi-cgbn}).
The 18 primes of 90--180 digits in \code{bench/primes.txt} ($4 \times 90$,
$5 \times 100$, $2 \times 110$, $1 \times 120$, $2 \times 130$, $2 \times
140$, $1 \times 150$, $1 \times 180$). Layouts of workers $\times$ ECM
threads $1 \times 8$, $4 \times 8$ and $8 \times 4$ for seeds 1--5, each
host-only and with \code{-C}, plus $1 \times 16$ and $1 \times 32$ for seed
1. Rank 0 coordinates; all ranks are unbound (\code{mpirun -{}-oversubscribe
-{}-bind-to none}), and every \code{-C} rank shares the one GPU. Host-only
and \code{-C} runs alternate, so slow drift affects both alike. Every run
proved all 18 numbers.

\begin{table}[ht]
\centering
\begin{tabular}{lrrrr}
\toprule
Workers $\times$ threads & Seeds & Host only & \code{-C} (CGBN) & Ratio \\
\midrule
$1 \times 8$  & 5 & 47.5\,s (39.3--57.3) & 25.9\,s (23.7--27.7) & $1.83\times$ \\
$1 \times 16$ & 1 & 33.1\,s & 24.2\,s & $1.37\times$ \\
$1 \times 32$ & 1 & 27.6\,s & 23.1\,s & $1.20\times$ \\
$4 \times 8$  & 5 & 17.0\,s (13.4--19.8) & 17.4\,s (15.7--18.6) & $0.98\times$ \\
$8 \times 4$  & 5 & 21.5\,s (20.1--24.9) & 18.5\,s (17.0--19.5) & $1.16\times$ \\
\midrule
$1 \times 8$, bound$^{*}$ & 1 & 185.3\,s & 51.7\,s & $3.59\times^{*}$ \\
\bottomrule
\end{tabular}
\caption{Batch wall time, mean (minimum--maximum) over the seeds.
  $^{*}$The worker bound to one core, Open MPI's default for two processes:
  an artifact (\S\ref{sec:binding}).}
\label{tab:mpi}
\end{table}

With several workers, the summed proof time is 52.7\,s host-only against
46.3\,s with CGBN at $4 \times 8$ ($1.14\times$), and 79.2\,s against
79.0\,s at $8 \times 4$ ($1.00\times$).

\begin{table}[ht]
\centering
\small
\begin{tabular}{rrrrr}
\toprule
\# & Digits & Host only & \code{-C} (CGBN) & Host/CGBN \\
\midrule
0  & 90  & 0.47 (0.38--0.50)   & 0.51 (0.49--0.57) & $0.92\times$ \\
1  & 90  & 0.69 (0.65--0.85)   & 1.46 (1.40--1.58) & $0.47\times$ \\
2  & 90  & 2.79 (2.68--2.97)   & 0.37 (0.36--0.38) & $7.54\times$ \\
3  & 90  & 0.64 (0.55--0.86)   & 0.38 (0.36--0.40) & $1.68\times$ \\
4  & 100 & 4.32 (4.19--4.49)   & 1.59 (0.60--1.94) & $2.72\times$ \\
5  & 100 & 0.25 (0.23--0.31)   & 0.45 (0.41--0.51) & $0.56\times$ \\
6  & 100 & 1.09 (0.50--5.80)   & 0.46 (0.29--0.50) & $2.37\times$ \\
7  & 100 & 0.36 (0.34--0.74)   & 0.70 (0.60--0.75) & $0.51\times$ \\
8  & 100 & 1.41 (1.13--1.60)   & 0.34 (0.29--0.45) & $4.15\times$ \\
9  & 110 & 1.44 (1.39--1.51)   & 0.93 (0.75--1.00) & $1.55\times$ \\
10 & 110 & 1.04 (0.89--1.31)   & 0.55 (0.44--0.71) & $1.89\times$ \\
11 & 120 & 1.09 (0.82--1.52)   & 1.37 (1.29--1.50) & $0.80\times$ \\
12 & 130 & 1.41 (1.33--5.56)   & 2.89 (1.49--2.93) & $0.49\times$ \\
13 & 130 & 1.19 (1.17--1.23)   & 2.55 (0.95--2.74) & $0.47\times$ \\
14 & 140 & 10.60 (7.56--13.12) & 1.38 (1.09--1.44) & $7.68\times$ \\
15 & 140 & 2.75 (1.46--5.86)   & 1.05 (0.97--1.16) & $2.62\times$ \\
16 & 150 & 7.46 (5.03--7.76)   & 1.43 (1.36--1.48) & $5.22\times$ \\
17 & 180 & 6.09 (3.42--9.26)   & 8.50 (7.12--9.33) & $0.72\times$ \\
\midrule
   &     & 45.1 (sum of medians) & 26.9           & $1.68\times$ \\
\bottomrule
\end{tabular}
\caption{Per proof, $1 \times 8$: median (minimum--maximum) over 5 seeds,
  seconds.}
\label{tab:per}
\end{table}

\subsection{What it shows}

\paragraph{The GPU's value falls as the host's thread budget rises:}
$1.83\times$, $1.37\times$ and $1.20\times$ with 8, 16 and 32 ECM threads
on one worker. The prover escalates to the GPU only when the host attempt
fails on a cofactor of 192 bits or more. With more threads, the host
attempt, which ends at the first factor any thread finds, succeeds more
often and sooner. The CGBN side barely moves ($25.9 \to 24.2 \to
23.1$\,s), because its time is dominated by work the extra host threads do
not touch.

\paragraph{Several ranks on one GPU dilute it to nothing in wall time:}
$0.98\times$ at $4 \times 8$ and $1.16\times$ at $8 \times 4$. The ranks
contend for one device, and each pays its own CUDA context creation and
module loading. Proof \#1 (90 digits) with \code{-C} takes 1.40--1.58\,s on
one worker, 3.23--3.58\,s with four and 6.46--7.17\,s with eight, across all
five seeds.

\paragraph{It works on the tail.} CGBN wins the hard proofs, those where
the host would grind on a stubborn cofactor: \#2 $7.5\times$, \#14
$7.7\times$, \#16 $5.2\times$, \#8 $4.2\times$. It loses the easy ones, 8 of
18, by $0.47$--$0.92\times$: each escalation pays fixed GPU costs
(launches, a 4096-curve wave, transfers) that the host would not have. The
result is variance reduction: across seeds the host-only batch ranged
39.3--57.3\,s, the CGBN batch 23.7--27.7\,s; per proof, \#6 ranged
0.50--5.80\,s on the host and 0.29--0.50\,s with CGBN. This is also what the
project's earlier single-thread sweep showed (\code{make bench-cuda}, 36
proofs, one ECM thread: $420.6 \to 120.8$\,s, $3.48\times$), in the
configuration most favourable to the GPU.

\subsection{Determinism, and why runs vary}

\code{gmpecpp\_prove()} reseeds its generator from the configured seed on
entry, so a proof does not depend on which rank takes it or in what order;
a seed fixes each proof's random choices (discriminant probing, points,
curves). With \code{-T > 1}, though, the host ECM is a race. Each thread
draws curves from its own deterministic stream, seeded from the proof seed
and the thread id, and the first thread to find a factor wins, under a
mutex and an atomic flag (\code{std::atomic<int> found}). Which factor wins
depends on thread timing, and a different factor of a candidate order
changes which candidates are accepted, and so the whole downrun. Under
contention, from several workers or anything else competing for cores, the
races can resolve differently from run to run; the seed ranges in
Table~\ref{tab:mpi} therefore include timing-induced variation as well as
the seeds' own. The certificates remain valid either way:
\code{gmpecpp-verify} checks them, where byte comparison cannot.

\subsection{The binding trap}
\label{sec:binding}

Open MPI binds processes to cores by default for small jobs. With two
processes (coordinator plus one worker):
\begin{lstlisting}
$ mpirun --report-bindings -n 2 ./gmpecpp-mpi -T 8 primes.txt
... Rank 0 bound to package[0][core:0]
... Rank 1 bound to package[0][core:1]
$ mpirun --oversubscribe --report-bindings -n 2 ./gmpecpp-mpi -T 8 primes.txt
... Rank 0 is not bound (or bound to all available processors)
... Rank 1 is not bound (or bound to all available processors)
\end{lstlisting}
The worker's threads inherit its CPU affinity: with \code{-T 8}, all eight
ECM threads time-share core 1's two hardware threads.

\begin{center}
\begin{tabular}{lrrr}
\toprule
$1 \times 8$, seed 1 & Host only & \code{-C} (CGBN) & Apparent ratio \\
\midrule
Unbound           & 45.1\,s  & 26.0\,s & $1.73\times$ \\
Bound to one core & 185.3\,s & 51.7\,s & $3.59\times$ \\
\bottomrule
\end{tabular}
\end{center}

The host path, nearly all threaded ECM, slows $4.1\times$; the CGBN path,
which needs little host CPU, slows $2\times$. The comparison then reports a
$3.6\times$ GPU advantage that is mostly lost host parallelism. The trap is
easy to fall into, because \code{-{}-oversubscribe} disables binding as a
side effect, so multi-rank jobs launched with it look fine. The harness
passes \code{-{}-bind-to none} explicitly. For multi-threaded MPI ranks
generally, check \code{-{}-report-bindings} (or \code{taskset -p} on a
worker) before trusting any thread-scaling number.

\section{Recommendations}

\begin{itemize}
  \item \textbf{Do not port scalar GMP code to CGBN.} Experiment 1 bounds it
    at $4\times$ to $172\times$ slower per operation.
  \item \textbf{Port only batched, same-width integer work.} ECM already is.
    Further candidates would have to be restructured into batches first
    (for example, testing many candidate orders or discriminants at once),
    and even then Amdahl's law applies to what remains.
  \item \textbf{For batch proving on one GPU,} prefer few workers with many
    unbound ECM threads each, and \code{-C}: CGBN then removes the
    pathological proofs. With many ranks, give each GPU one rank
    (\code{gmpecpp-mpi} across nodes with one GPU each is the natural
    layout).
  \item \textbf{The factorizer gains nothing} from moving its remaining host
    GMP (at most 0.06\%).
\end{itemize}

\section{What was wrong along the way}

Five intermediate results in this investigation were wrong. Each was caught
by measuring again with a controlled setup, rather than by reasoning about
the first numbers.
\begin{enumerate}
  \item \textbf{``CGBN is 3.6--4$\times$ faster with one worker.''} The
    first single-worker runs used \code{mpirun -n 2} without
    \code{-{}-oversubscribe}, so the worker was bound to one core
    (\S\ref{sec:binding}).
  \item \textbf{``\code{-T 32} is slower than \code{-T 8} because of
    E-cores and \code{curves\_per\_bound}.''} The explanation was plausible
    --- $\max(T, 20)$ curves per bound adds work above 20 threads, and
    threads beyond 8 land on SMT siblings and E-cores --- but its premise
    was the binding artifact. Unbound, 32 threads beat 8 (27.6\,s against
    45.1\,s).
  \item \textbf{``Multi-rank runs differ because the numbers reach ranks in
    a different order and the random state carries over.''} The prover
    reseeds at every proof. The real source is the threaded ECM race under
    contention.
  \item \textbf{``With \code{-C}, the GMP that remains is the class
    polynomials' \code{mpf}.''} This came from misreading an exploratory
    \code{perf} report. With \code{cpu-clock} and the symbols grouped,
    \code{mpf} is 0.1\%, and gcd and division from the affine point
    arithmetic is 52\%.
  \item \textbf{Profile shares that added up to more than 100\%.} On the
    hybrid CPU, \code{perf}'s default \code{cycles} event produced two report
    tables (P-core and E-core PMUs), and the first grouping summed across
    both. \code{cpu-clock} fixed it.
\end{enumerate}

\section{Reproducing}

\begin{lstlisting}
make                       # libraries, gmpecpp-mpi (needs Open MPI), gmpecmfft
make bench-cgbn-latency    # experiment 1, ~1 min
make bench-gmp-profile     # experiment 2, ~2 min; needs perf
make bench-mpi-cgbn        # experiment 3, ~40 min; build/bench-mpi-cgbn/
bash bench/mpi-cgbn.sh DIR "1 2 3" bound   # parts and seeds can be chosen
\end{lstlisting}
\code{bench-mpi-cgbn} writes each run's raw output to
\path{build/bench-mpi-cgbn/<layout>-<cpu|gpu>-s<seed>.txt} and the
tables to \code{summary.txt}. Numbers will differ on other hardware; the ratios
and the binding effect are what to compare.

\end{document}
